% =====================================================================
%  9. Calendarización   (actividad TP6)
% =====================================================================
\section{Calendarización}
\label{sec:calendario}

\begin{instruccion}
Elabore el cronograma a partir de las actividades y duraciones del
Cuadro~\ref{tab:actividades}, de sus dependencias y de la disponibilidad declarada en el
apartado~\ref{sec:equipo}. Incluya revisiones, correcciones, entregas
y evaluaciones, así como los hitos del Cuadro~\ref{tab:hitos}.

Puede sustituir el cuadro siguiente por un diagrama de Gantt, siempre que éste lleve caption,
etiqueta y una mención desde el texto. Si lo incluye como imagen, colóquela en
\texttt{diagrams/} o en una carpeta \texttt{img/} y refiérala con ruta relativa.
\end{instruccion}

\begin{table}[H]
\centering
\caption{Cronograma planificado del esfuerzo de prueba.}
\label{tab:cronograma}
\begin{tabularx}{\textwidth}{|C{1.4cm}|Y|C{2.5cm}|C{2.5cm}|C{2.4cm}|}
  \hline
  \filaenc \textbf{Clave} & \textbf{Actividad o hito} & \textbf{Inicio} & \textbf{Fin} & \textbf{Responsable} \\ \hline
  \campo{ACT-01} & \campo{Actividad} & \campo{dd/mm/aaaa} & \campo{dd/mm/aaaa} & \campo{Rol} \\ \hline
  \campo{H-01}   & \campo{Primera entrega} & \campo{dd/mm/aaaa} & \campo{dd/mm/aaaa} & \campo{Rol} \\ \hline
  & & & & \\ \hline
  & & & & \\ \hline
\end{tabularx}
\end{table}

El Cuadro~\ref{tab:cronograma} sitúa en el tiempo las actividades y los hitos previstos, y
constituye la línea base temporal contra la que el apartado~\ref{sec:seguimiento} compara el
avance real. Al cierre, se contrasta con las fechas efectivamente ejecutadas en el
apartado~\ref{sec:resumen-final}.
